\section{Related Works}
\label{sec:related_works}

\paragraph{Safety Degradation under Fine-Tuning.}
\citet{qi2024finetuning} show that fine-tuning on benign data compromises safety. \citet{qi2025tokensdeep} find alignment is shallow, encoded in only a few token positions. Defenses include LISA~\cite{huang2024lisa} (lazy alignment during training), Safe LoRA~\cite{hsu2024safe} (projection of LoRA updates), and Vaccine~\cite{yi2024vaccine} (perturbation-aware alignment). All require per-checkpoint intervention.

\paragraph{Safety Classifiers and Guardrails.}
Llama Guard~\cite{inan2023llama} and WildGuard~\cite{han2024wildguard} classify prompts independently of the target model.
HiddenDetect~\cite{jiang2025hiddendetect} monitors hidden states but requires per-model probe fitting.
These methods are either model-agnostic (ignoring fine-tuning effects) or per-model (not transferable).

\paragraph{Hypernetworks and Side Networks.}
Hypernetworks~\cite{ha2017hypernetworks, chauhan2024brief} generate target-network weights conditioned on context, with applications in continual learning~\cite{vonoswald2020continual} and multi-task adaptation~\cite{navon2023equivariant}. Recent work has scaled this paradigm to LLMs: \citet{liang2025dragdrop} generate LoRA parameters directly from task descriptions, achieving zero-shot adaptation without gradient-based fine-tuning. Our work pursues a complementary direction---rather than generating task-adaptation weights from text prompts, we generate \textit{safety} weights from model activations, targeting the orthogonal problem of safety restoration.
Ladder Side-Tuning (LST)~\cite{sung2022lst} attaches a small side network to a frozen backbone for parameter-efficient transfer. We are, to our knowledge, the first to combine hypernetworks with side-network architectures for safety restoration.